\section{矩阵与逆}
\idx{矩阵}{juzhen}
\idx{逆矩阵}{nijuzhen}

把矢量写成列.矩阵 $M$ 作用在 $\mathbf{x}$ 上,是各列的线性组合,
\begin{equation}
  (M\mathbf{x})_i=M_{ij}x_j.
\end{equation}
重复指标 $j$ 要求和,这是下一节的求和约定,这里先按这个意思使用.相继两次映射先做 $N$ 再做 $M$,
\begin{equation}
  M(N\mathbf{x})=(MN)\mathbf{x}.
\end{equation}
单位矩阵的列就是基矢,$I\mathbf{x}=\mathbf{x}$.单位正方形被 $M$ 送到一个平行四边形,面积是 $|\det M|$.

\subsection{逆矩阵}
若存在 $A^{-1}$ 使 $A A^{-1}=A^{-1}A=I$,就说 $A$ 可逆,线性映射可以做回去.逆存在当且仅当 $\det A\neq 0$.用余子式构造:划去第 $i$ 行第 $j$ 列,剩下的行列式乘上 $(-1)^{i+j}$,得到余子式 $C_{ij}$.伴随矩阵是 $\operatorname{adj} A=C^{\mathsf T}$,并且
\begin{equation}
  A\,(\operatorname{adj} A)=(\det A)\,I,
  \qquad
  A^{-1}=\frac{1}{\det A}\operatorname{adj} A.
\end{equation}

$2\times 2$ 可以直接验证.设
\begin{equation}
  A=\begin{pmatrix} a & b \\ c & d \end{pmatrix},
  \qquad
  \det A=ad-bc.
\end{equation}
\begin{equation}
  \begin{pmatrix} a & b \\ c & d \end{pmatrix}
  \begin{pmatrix} d & -b \\ -c & a \end{pmatrix}
  =(ad-bc)I,
\end{equation}
所以
\begin{equation}
  A^{-1}=\frac{1}{ad-bc}\begin{pmatrix} d & -b \\ -c & a \end{pmatrix}.
\end{equation}

\begin{example}
$A=\begin{pmatrix}2&1\\1&2\end{pmatrix}$,$\mathbf{b}=\begin{pmatrix}1\\0\end{pmatrix}$.
\end{example}
\begin{solution}
$\det A=3$,
\begin{equation}
  A^{-1}=\frac13\begin{pmatrix}2&-1\\-1&2\end{pmatrix}.
\end{equation}
直接相乘,$AA^{-1}=I$.方程 $A\mathbf{x}=\mathbf{b}$ 的解是
\begin{equation}
  \mathbf{x}=A^{-1}\mathbf{b}=\begin{pmatrix}2/3\\-1/3\end{pmatrix}.
\end{equation}
\end{solution}

\begin{example}
\begin{equation}
  A=\begin{pmatrix} 2&1&0\\ 1&2&1\\ 0&1&2 \end{pmatrix}.
\end{equation}
\end{example}
\begin{solution}
沿第一行展开,$\det A=2(4-1)-(2-0)=4$.余子式矩阵是
\begin{equation}
  C=\begin{pmatrix} 3&-2&1\\ -2&4&-2\\ 1&-2&3 \end{pmatrix}.
\end{equation}
它已经对称,$\operatorname{adj} A=C$,
\begin{equation}
  A^{-1}=\frac14 C.
\end{equation}
$A C$ 的 $(1,1)$ 元是 $2\cdot 3+1\cdot(-2)=4=\det A$,$(1,2)$ 元是 $2\cdot(-2)+1\cdot 4=0$.这与 $A\,(\operatorname{adj} A)=(\det A)\,I$ 一致.
\end{solution}

若 $R^{\mathsf T}R=I$,则 $R^{-1}=R^{\mathsf T}$,不必再做余子式.旋转矩阵
\begin{equation}
  R=\begin{pmatrix} \cos\theta & -\sin\theta \\ \sin\theta & \cos\theta \end{pmatrix}
\end{equation}
就是这种正交矩阵.它保持矢量的模和夹角,核对放在指标记号里.
